\textbf{整体二分}

整体二分是一种离线分治思想。

\textbf{适用条件}

\[
\text{答案可二分，判定有单调性，并且题目允许离线。}
\]

常见判定形式是：对当前 \(mid\)，判断每个询问的答案是否 \(\le mid\)。

\textbf{递归形式}

设当前处理状态为
\[
F(L,R,Q),
\]
其中 \([L,R]\) 是答案值域，\(Q\) 是当前操作与询问序列。

若 \(L=R\)，则 \(Q\) 中所有询问答案都是 \(L\)。

否则令 \(mid=\lfloor(L+R)/2\rfloor\)，统计所有操作对“答案 \(\le mid\)”的贡献，
并把 \(Q\) 稳定划分为
\[
Q_L,\ Q_R,
\]
其中 \(Q_L\) 的答案在 \([L,mid]\)，\(Q_R\) 的答案在 \([mid+1,R]\)。
然后递归处理
\[
F(L,mid,Q_L),\qquad F(mid+1,R,Q_R).
\]

\textbf{复杂度}

若每层分流与统计的代价为 \(T(n)\)，值域大小为 \(V\)，
则总复杂度通常为
\[
O(T(n)\log V).
\]

常见情形是每层用树状数组或线段树线性扫一遍操作序列，
于是复杂度常写作
\[
O(n\log n\log V)
\]
或
\[
O((n+m)\log n\log V).
\]
